\section{Cơ sở lý thuyết}
\subsection{Mô hình thương mại điện tử B2C, B2B, B2B2C và marketplace nhiều người bán}


Trong thương mại điện tử, các mô hình được phân loại chủ yếu theo bản chất của các bên tham gia giao dịch.. B2B (Business to Business) là mô hình kinh doanh giữa các doanh nghiệp với nhau. Trong đó doanh nghiệp cung cấp sản phẩm hoặc dịch vụ cho các doanh nghiệp khác thay vì cho người tiêu dùng cuối. Đặc trưng của mô hình này là ở các đơn hàng có số lượng lớn và có tính chuyên nghiệp cao thông qua hợp đồng. Trái lại, B2C (Business to Consumer) là mô hình kinh doanh giữa doanh nghiệp và người tiêu dùng cá nhân. Các công ty cung cấp sản phẩm hoặc dịch vụ trực tiếp cho người tiêu dùng cuối cùng như các cửa hàng bán lẻ. B2B2C (Business to Business to Consumer) là một mô hình lai (hybrid), kết hợp giữa B2B và B2C trong đó một doanh nghiệp cung cấp sản phẩm hoặc dịch vụ cho một doanh nghiệp khác, và doanh nghiệp thứ hai này sẽ bán trực tiếp cho người tiêu dùng. Đây là mô hình phổ biến trong các nền tảng thương mại điện tử như các sàn giao dịch trực tuyến, nơi các bên thứ ba có thể bán sản phẩm cho người tiêu dùng. Ngoài ra, Multi vendor marketplace (sàn nhiều nhà cung cấp) là mô hình thương mại điện tử cho phép nhiều người bán (vendor) đăng ký, đăng sản phẩm/dịch vụ, quản lý đơn hàng và nhận thanh toán trên cùng một nền tảng. Người sở hữu sàn (marketplace operator) đóng vai trò quản trị, định ra quy tắc, thu phí và đảm bảo trải nghiệm mua bán liền mạch.


Farmery lựa chọn mô hình B2B2C dạng marketplace nhiều người bán nhằm đáp ứng nhu cầu mua nông sản của 2 tệp khách hàng khác nhau: người tiêu dùng cá nhân mua lẻ và doanh nghiệp mua sỉ theo hợp đồng từ ngồn cung là nông dân/ hợp tác xã, còn nền tảng sẽ đóng vai trò trung gian cung cấp công nghệ, là nền tảng thanh toán và là nơi đáng tin cậy để đảm bảo các giao dịch diễn ra đúng quy trình, mang lại lợi ích cho các bên.


\subsection{Lý thuyết chuỗi cung ứng nông sản ngắn (SFSC) và disintermediation}


Chuỗi cung ứng thực phẩm ngắn (SFSC), theo định nghĩa của EU, là một chuỗi cung ứng liên quan đến một số lượng hạn chế các nhà điều hành kinh tế, cam kết hợp tác, phát triển kinh tế địa phương và duy trì mối quan hệ gần gũi về địa lý và xã hội giữa các nhà sản xuất thực phẩm, nhà chế biến và người tiêu dùng \cite{tapchicongsan2018}. Có một số hình thức khác nhau của SFSC. Một trong những hình thức đơn giản nhất là bán trực tiếp từ nông dân đến người tiêu dùng cuối cùng (tại nông trại, chợ nông dân, giao hàng qua internet). Các hình thức khác bao gồm các chương trình giao hộp sản phẩm, "tự hái", và nông nghiệp được cộng đồng hỗ trợ (CSA), nơi người tiêu dùng hỗ trợ tài chính cho các nhà trồng trọt địa phương bằng cách mua "gói đăng ký" các sản phẩm tươi cho một mùa vụ nhất định. Các sản phẩm chính thường được giao dịch trong SFSC là trái cây và rau củ theo mùa, tiếp theo là sản phẩm từ động vật (chủ yếu là thịt, tươi và đã chế biến) và các sản phẩm từ sữa cũng như đồ uống. Lợi ích chính của SFSC là tăng tính minh bạch và niềm tin cho người tiêu dùng, đồng thời mang lại lợi nhuận kinh tế tốt hơn, mang lại quyền tự chủ lớn hơn cho nhà sản xuất đồng thời tăng cường khả năng lưu thông của hàng hóa; hạn chế chính là khối lượng sản xuất/phân phối bị giới hạn, chi phí vận hành trên mỗi đơn vị sản phẩm cao hơn và khó mở rộng quy mô so với chuỗi cung ứng dài truyền thống.


Một khái niệm liên quan đến SFSC là disintermediation (loại bỏ trung gian) — quá trình mà công nghệ số cho phép nhà sản xuất tiếp cận trực tiếp người mua cuối, cắt giảm các tầng trung gian truyền thống (thương lái, chợ đầu mối, nhà phân phối cấp trung). Một ví dụ điển hình về việc loại bỏ trung gian là chiến lược mà nhà bán lẻ máy tính trực tuyến Dell áp dụng vào đầu thế kỷ 21 \cite{britannica_disintermediation}. Công ty bán hàng hóa qua trang web của mình nhưng không có mặt trực tiếp tại các trung tâm mua sắm. Việc tiết kiệm chi phí vận hành đã giúp công ty cung cấp nhiều loại hàng hóa với giá thấp hơn so với các nhà bán lẻ truyền thống. Farmery về bản chất là một nỗ lực disintermediation có kiểm soát: không loại bỏ hoàn toàn trung gian, vẫn là nền tảng hợp tác với các đơn vị vận chuyển, tuy nhiên nền tảng sẽ là nơi đứng ra là nơi trung gian cung cấp công nghệ, đảm bảo uy tín để mang lại sự yên tâm cho người dùng. Điều này giúp loại bỏ các tầng trung gian thương mại thuần túy hưởng chênh lệch giá mà không tạo thêm giá trị tương xứng như thương lái ép giá và vẫn giữ được một số lợi ích quy mô của chuỗi cung ứng dài (tiếp cận thị trường rộng hơn phạm vi địa lý của một chuỗi ngắn truyền thống) nhờ nền tảng số.

\subsection{Truy xuất nguồn gốc thực phẩm: chuẩn GS1 và vai trò của blockchain}


GS1 là tổ chức quốc tế phi lợi nhuận xây dựng các chuẩn định danh và trao đổi dữ liệu dùng chung toàn cầu cho chuỗi cung ứng, trong đó có bộ chuẩn dành riêng cho truy xuất nguồn gốc (GS1 Traceability Standard). GS1 định nghĩa cả mã định danh và phương tiện lưu trữ dữ liệu . Mã định danh bao gồm GTIN cho sản phẩm ở mọi cấp độ đóng gói, GLN cho các bên và địa điểm, SSCC cho các đơn vị hậu cần và các khóa chuyên dụng cho tài sản và mối quan hệ. Sau khi có định danh, hệ thống sẽ ghi lại thông tin chuỗi cung ứng gắn với các định danh này, cho phép truy xuất thông tin lô hàng khi cần thu hồi hoặc xác minh xuất xứ \cite{gs1_keys}. Farmery ứng dụng điều này và tạo ra các mã QR gắn với từng lô hàng, đại diện cho khả năng định danh nguồn gốc, ghi lại dữ liệu ứng với lô hàng đó.


Blockchain là một công nghệ chuỗi khối phân tán, cho phép ghi chép và lưu trữ thông tin theo cách không thể thay đổi hoặc làm giả. Nhờ vậy, các giao dịch được thực hiện trên nền tảng blockchain có thể đảm bảo tính minh bạch, tin cậy và bảo mật cao. Trong lĩnh vực thương mại điện tử, blockchain có thể hỗ trợ các bên (doanh nghiệp, nông dân, đơn vị vận chuyển) trong việc truy xuất nguồn gốc sản phẩm, tối ưu hóa quản lý chuỗi cung ứng và nâng cao hiệu quả giao dịch thương mại quốc tế. Với sự phát triển mạnh mẽ của thương mại điện tử và sự gia tăng của các giao dịch xuyên biên giới, blockchain có thể đóng vai trò như một công cụ giúp giảm thiểu gian lận, bảo vệ quyền lợi của người tiêu dùng và cải thiện môi trường kinh doanh.


Tuy nhiên, các nghiên cứu gần đây chỉ ra rằng việc áp dụng blockchain cho truy xuất nguồn gốc thực phẩm vẫn gặp rào cản thực tế: Bài toán "Garbage In, Garbage Out" và sai số tại ranh giới vật lý – kỹ thuật số. Blockchain chỉ đảm bảo dữ liệu không thể bị chỉnh sửa sau khi đã ghi, chứ không thể tự động xác minh dữ liệu đưa vào ban đầu có chính xác hay không (input verification). Trong chuỗi cung ứng thực phẩm, nông dân hoặc các cơ sở chế biến nhỏ lẻ vẫn phải nhập thông tin thủ công (như loại phân bón sử dụng, ngày thu hoạch, chứng nhận hữu cơ) hoặc gắn thẻ RFID/QR code lên sản phẩm. Hiện tượng gian lận tại nguồn, dán nhầm nhãn, hoặc cố tình tráo đổi hàng kém chất lượng trước khi đóng gói vẫn diễn ra mà blockchain không phát hiện được \cite{bath_blockchain}. Với quy mô một sàn TMĐT nông sản mới như Farmery, việc lưu trữ nhật ký canh tác và dữ liệu truy xuất trên một cơ sở dữ liệu tập trung tuân theo nguyên lý định danh của GS1 là phương án khả thi và tiết kiệm chi phí hơn ở giai đoạn đầu; blockchain có thể được xem xét bổ sung ở giai đoạn sau nếu cần tăng độ tin cậy cho các đối tác xuất khẩu hoặc khi số lượng bên tham gia độc lập vào chuỗi tăng đủ lớn để việc phi tập trung hóa mang lại lợi ích rõ ràng, vượt qua chi phí vận hành bổ sung.

\subsection{Hệ thống đánh giá, uy tín và niềm tin (reputation, information asymmetry)}


Nền tảng lý thuyết cho vấn đề đảm bảo có một hệ thống đánh giá đủ uy tín được xây dựng dựa trên khái niệm bất cân xứng thông tin (information asymmetry), nổi bật với mô hình kinh điển "The Market for Lemons" của Akerlof (1970). Theo đó, khi người bán nắm rõ chất lượng thực tế của sản phẩm hơn người mua (ví dụ: mức độ tuân thủ tiêu chuẩn canh tác an toàn) và người mua thiếu công cụ để kiểm chứng trước khi giao dịch, thị trường sẽ đối mặt với rủi ro "lựa chọn đối nghịch" (adverse selection) – nơi hàng kém chất lượng dần đào thải hàng hóa tốt. Nguyên nhân là do những người bán sản phẩm chất lượng cao không thể chứng minh giá trị thực để thuyết phục mức giá tương xứng, buộc họ phải rời khỏi thị trường hoặc hạ chuẩn chất lượng để cạnh tranh về giá. Đối với một sàn thương mại điện tử nông sản như Farmery, đây là bài toán cốt lõi cần giải quyết: nếu thiếu vắng các cơ chế minh bạch hóa, người tiêu dùng sẽ hoàn toàn mất phương hướng trong việc phân biệt nông sản "sạch" thực sự với các sản phẩm có mục đích marketing.


Hệ thống đánh giá và uy tín (reputation system) là cơ chế phổ biến để giảm bất cân xứng thông tin trên các sàn giao dịch trực tuyến, bằng cách tích lũy phản hồi từ các giao dịch trước đó (điểm số, bình luận, tỷ lệ hoàn thành đơn) thành một chỉ số uy tín công khai cho từng người bán. Nghiên cứu tổng quan của Tadelis về các hệ thống danh tiếng và phản hồi trên các nền tảng trực tuyến chỉ ra rằng các cơ chế này giúp thị trường vận hành hiệu quả hơn bằng cách tạo động lực cho người bán duy trì chất lượng để bảo vệ uy tín dài hạn, dù vẫn tồn tại vấn đề như tâm lý ngại đánh giá tiêu cực do sợ bị trả đũa, hoặc đánh giá tập trung ở nhóm khách hàng cực đoan (rất hài lòng hoặc rất không hài lòng). Với Farmery, hệ thống đánh giá cần kết hợp với dữ liệu truy xuất nguồn gốc khách quan và chứng nhận đã được xác minh như hai lớp bổ sung cho đánh giá chủ quan của người dùng, giúp giảm thiểu rủi ro đánh giá giả hoặc thao túng, đặc biệt quan trọng với nhóm khách B2B khi giá trị mỗi giao dịch lớn hơn nhiều so với B2C.


Hệ thống đánh giá và uy tín (reputation system) là cơ chế phổ biến để giảm bất cân xứng thông tin trên các sàn giao dịch trực tuyến, bằng cách tích lũy phản hồi từ các giao dịch trước đó (điểm số, bình luận, tỷ lệ hoàn thành đơn) thành một chỉ số uy tín công khai cho từng người bán \cite{springer_article}. Với Farmery, hệ thống đánh giá cần kết hợp với dữ liệu truy xuất nguồn gốc khách quan và chứng nhận đã được xác minh như hai lớp bổ sung cho đánh giá chủ quan của người dùng, giúp giảm thiểu rủi ro đánh giá giả hoặc thao túng, đặc biệt quan trọng với nhóm khách B2B khi nhóm khác này thường giao dịch với số lượng lớn và giá trị các giao dịch cao hơn so với B2C.

\subsection{Nền tảng công nghệ: kiến trúc web, RBAC, cổng thanh toán trung gian và escrow}

Về kiến trúc web, với một sàn thương mại điện tử đa vai trò như Farmery, để có thể phục vụ nhiều nhóm người dùng gồm nông dân/HTX, người tiêu dùng, doanh nghiệp B2B và quản trị viên, có thể được xây dựng theo mô hình kiến trúc phân lớp hoặc kiến trúc hướng dịch vụ, tùy thuộc vào quy mô và yêu cầu mở rộng của hệ thống. Các thành phần chính bao gồm: lớp giao diện cung cấp web/app tương ứng với từng nhóm người dùng; lớp API và nghiệp vụ chịu trách nhiệm xử lý các quy trình kinh doanh; lớp dữ liệu sử dụng cơ sở dữ liệu quan hệ để lưu trữ thông tin người dùng, sản phẩm, đơn hàng, giao dịch và hợp đồng, kết hợp với hệ thống lưu trữ đối tượng cho hình ảnh, tài liệu và các chứng nhận liên quan. Ngoài ra, hệ thống có thể tích hợp với các dịch vụ bên thứ ba như cổng thanh toán, đơn vị vận chuyển và các hệ thống xác thực hoặc kiểm tra chứng nhận thông qua API khi cần thiết.


Đối với kiểm soát truy cập, Farmery áp dụng mô hình Role-Based Access Control (RBAC). Đây là mô hình trong đó quyền truy cập không được cấp trực tiếp cho từng người dùng mà được tổ chức thông qua các vai trò (role). Người dùng được gán một hoặc nhiều vai trò, trong khi mỗi vai trò được liên kết với một tập hợp các quyền tương ứng với chức năng nghiệp vụ. Mô hình RBAC được Sandhu và cộng sự hệ thống hóa thành một khung lý thuyết vào năm 1996. Với đặc thù của Farmery, RBAC phù hợp để quản lý tối thiểu bốn nhóm vai trò chính: nông dân/HTX, người tiêu dùng, doanh nghiệp thu mua và quản trị viên, tương ứng với các nhóm người dùng đã được phân tích. Khi hệ thống phát triển, có thể tiếp tục phân tách các vai trò thành những vai trò chuyên biệt hơn, chẳng hạn quản trị viên kiểm duyệt nội dung và quản trị viên xử lý vi phạm. Cách tổ chức này giúp áp dụng nguyên tắc đặc quyền tối thiểu (Principle of Least Privilege), theo đó mỗi vai trò chỉ được cấp những quyền cần thiết để thực hiện nhiệm vụ, qua đó hạn chế nguy cơ truy cập trái phép, sai sót trong vận hành và các rủi ro bảo mật.


Về thanh toán, trung gian thanh toán (payment intermediary) là các tổ chức được cấp phép cung cấp dịch vụ hỗ trợ thu hộ/chi hộ và chuyển tiền giữa các bên trong giao dịch điện tử mà không phải là ngân hàng trực tiếp của cả hai bên (ví như các cổng ví điện tử VNPay, Momo tại Việt Nam). Bên cạnh hình thức thanh toán thông thường, Farmery có thể xem xét áp dụng cơ chế escrow (ký quỹ trung gian) cho các giao dịch có giá trị lớn, đặc biệt là giao dịch B2B. Với cơ chế này, khoản tiền của người mua được một bên trung gian giữ lại và chỉ được giải ngân cho người bán khi các điều kiện đã thỏa thuận được đáp ứng, chẳng hạn người mua xác nhận đã nhận đủ hàng và hàng hóa đạt yêu cầu về số lượng, chất lượng hoặc các tiêu chí được quy định trong hợp đồng. Cơ chế này giúp giảm rủi ro cho cả hai phía: doanh nghiệp thu mua hạn chế nguy cơ thanh toán nhưng không nhận được hàng hóa đúng cam kết, trong khi nông dân/HTX có thêm sự đảm bảo rằng khoản tiền đã được người mua nộp và được giữ bởi bên trung gian thay vì phụ thuộc hoàn toàn vào khả năng thanh toán sau này của doanh nghiệp.


\section*{Tài liệu trích dẫn:}


\nocite{*}


\bibliographystyle{ieeetr}


\bibliography{references}